\documentclass[../DoAn.tex]{subfiles}
\begin{document}

Chương này giới thiệu tổng quan về đề tài xây dựng hệ thống thương mại điện tử chuyên biệt cho thú cưng mang tên ThePawsome. Nội dung chương tập trung làm rõ tính cấp thiết của đề tài thông qua việc phân tích thực trạng thị trường và nhu cầu thực tế của người nuôi thú cưng tại Việt Nam. Trên cơ sở đó, chương xác định rõ ràng mục tiêu, phạm vi nghiên cứu, đề xuất định hướng giải pháp công nghệ cụ thể dựa trên nền tảng trí tuệ nhân tạo và các cơ chế giao dịch an toàn, đồng thời trình bày cấu trúc bố cục tổng thể của toàn bộ báo cáo đồ án tốt nghiệp.

\section{Đặt vấn đề}
\label{section:1.1}

Trong những năm gần đây, cùng với đà tăng trưởng kinh tế và xu hướng đô thị hóa nhanh chóng tại Việt Nam, nhu cầu nuôi dưỡng và chăm sóc thú cưng đã có sự bùng nổ mạnh mẽ. Thú cưng giờ đây không chỉ đơn thuần là vật nuôi giữ nhà, mà đã trở thành những thành viên quan trọng trong các gia đình, đòi hỏi sự quan tâm đặc biệt về sức khỏe và tinh thần. Sự thay đổi sâu sắc trong nhận thức của người nuôi đã thúc đẩy quy mô thị trường dịch vụ và sản phẩm dành cho thú cưng phát triển với tốc độ nhanh chóng. Theo đó, nhu cầu mua sắm các loại thức ăn dinh dưỡng chất lượng cao, phụ kiện chuyên dụng, dược phẩm thú y và các dịch vụ y tế ngày một tăng lên.

Tuy nhiên, người nuôi thú cưng tại Việt Nam đang phải đối mặt với nhiều khó khăn trong việc tiếp cận các nguồn lực chăm sóc chất lượng. Khó khăn lớn nhất nằm ở sự thiếu hụt thông tin tư vấn chính thống và cá nhân hóa. Mỗi cá thể thú cưng có những đặc điểm sinh lý, thể trạng, độ tuổi, thậm chí là tình trạng dị ứng thực phẩm rất khác nhau. Việc lựa chọn sản phẩm phù hợp nếu chỉ dựa vào cảm tính hoặc thông tin chung chung trên bao bì có thể dẫn đến những nguy cơ tiềm ẩn cho sức khỏe vật nuôi. Bên cạnh đó, các website thương mại điện tử chuyên biệt về thú cưng hiện nay tại thị trường trong nước hầu hết chỉ dừng lại ở mô hình mua bán truyền thống, thiếu các tính năng hỗ trợ thông minh, tương tác trực quan hoặc tư vấn chuyên sâu mang tính cá nhân hóa cho từng khách hàng.

Thêm vào đó, việc tìm kiếm và chia sẻ kiến thức chăm sóc thú cưng trên các mạng xã hội hay diễn đàn tự phát thường gặp phải vấn đề nhiễu loạn thông tin. Nhiều kiến thức y khoa hoặc chế độ dinh dưỡng không có sự kiểm chứng của những người có chuyên môn y tế thú y, dẫn đến việc áp dụng sai cách, ảnh hưởng tiêu cực tới sức khỏe vật nuôi. Người nuôi thú cưng luôn có nhu cầu trao đổi, kết nối trực tiếp với các chuyên gia có uy tín và những người có kinh nghiệm thực tế trong một môi trường thông tin được kiểm duyệt chặt chẽ.

Những vấn đề nêu trên nếu được giải quyết triệt để sẽ mang lại lợi ích to lớn cho cộng đồng người nuôi thú cưng tại Việt Nam. Việc hỗ trợ thông tin tư vấn cá nhân hóa và cảnh báo dị ứng tự động theo hồ sơ sức khỏe thú cưng giúp bảo vệ an toàn cho vật nuôi, đồng thời tiết kiệm thời gian và chi phí cho chủ nuôi. Một không gian trao đổi tri thức đáng tin cậy có sự xác thực của chuyên gia y tế sẽ góp phần nâng cao nhận thức cộng đồng, thúc đẩy xu hướng chăm sóc thú cưng một cách khoa học, an toàn và chuyên nghiệp.

\section{Mục tiêu và phạm vi đề tài}
\label{section:1.2}

Qua khảo sát, các sản phẩm thương mại điện tử petshop hiện có tại thị trường Việt Nam hầu hết chỉ tập trung vào chức năng mua bán trực tuyến đơn thuần. Mặc dù một số hệ thống có tích hợp các chuyên mục chia sẻ kiến thức cẩm nang, các bài viết này chủ yếu tồn tại dưới dạng thông tin tĩnh, thiếu khả năng tương tác trực tiếp và khó cá nhân hóa theo thể trạng của từng vật nuôi. Các công cụ tìm kiếm sản phẩm trên nhiều hệ thống hiện tại cũng thiên về từ khóa đơn giản, chưa khai thác tốt ngữ nghĩa câu hỏi hoặc ngữ cảnh cụ thể của khách hàng, gây ra nhiều khó khăn trong việc lựa chọn sản phẩm phù hợp với nhu cầu thực tế.

Dựa trên các phân tích và đánh giá ở trên, đồ án khái quát lại những hạn chế cốt lõi bao gồm sự thiếu hụt tính cá nhân hóa trong tư vấn sản phẩm, nguy cơ mua nhầm sản phẩm chứa thành phần dị ứng với thể trạng thú cưng, sự nhiễu loạn thông tin trên các diễn đàn tự phát và độ trễ lớn trong khâu hỗ trợ khách hàng khi gặp sự cố sức khỏe vật nuôi. Trên cơ sở đó, đồ án hướng tới mục tiêu thiết kế và xây dựng hệ thống thương mại điện tử chuyên biệt mang tên ThePawsome, tích hợp trợ lý ảo thông minh sử dụng kiến trúc RAG, agent tool-calling có planner và diễn đàn cộng đồng tương tác được kiểm chứng bởi chuyên gia y tế thú y. Hệ thống được phát triển với các nhóm chức năng chính bao gồm quản lý tài khoản và phân quyền, quản lý hồ sơ thú cưng chi tiết làm cơ sở cá nhân hóa, hệ thống danh mục sản phẩm biến thể phức tạp, quy trình giỏ hàng và đặt hàng an toàn kết hợp giữ kho tạm thời cùng thanh toán trực tuyến, diễn đàn thảo luận được xác thực và trợ lý ảo Catbot hỗ trợ thời gian thực có khả năng tự động chuyển giao cuộc trò chuyện sang tư vấn viên là con người.

Trong phần này, phạm vi nghiên cứu công nghệ của đồ án được giới hạn trong việc ứng dụng mô hình ngôn ngữ lớn kết hợp kiến trúc RAG dựa trên cơ sở dữ liệu vector PostgreSQL và phần mở rộng pgvector. Đồ án cũng tập trung nghiên cứu mô hình tìm kiếm lai Hybrid Search trộn điểm bằng thuật toán Reciprocal Rank Fusion, thiết lập hệ thống quản lý hội thoại đa tác vụ phức tạp sử dụng LangGraph, và áp dụng cơ chế bộ đệm Redis để tối ưu hóa hiệu năng phản hồi hệ thống dưới mức giới hạn cho phép.

\section{Định hướng giải pháp}
\label{section:1.3}

Từ việc xác định rõ nhiệm vụ cần giải quyết ở phần \ref{section:1.2}, đồ án đề xuất định hướng giải pháp của mình theo trình tự khoa học cụ thể. Trước tiên, đồ án lựa chọn giải quyết vấn đề theo định hướng kiến trúc ứng dụng web phân tách rõ ràng giữa lớp giao diện Frontend và lớp xử lý Backend thông qua hệ thống RESTful API, kết hợp với các kỹ thuật trí tuệ nhân tạo hiện đại.

Tiếp theo, giải pháp kỹ thuật cụ thể của hệ thống được mô tả ngắn gọn thông qua sự kết hợp đồng bộ của các công nghệ cốt lõi. Lớp Backend được xây dựng trên ngôn ngữ Python với framework FastAPI và hệ quản trị cơ sở dữ liệu PostgreSQL tích hợp tiện ích pgvector, kết hợp Redis làm bộ nhớ đệm. Lớp Frontend sử dụng framework Next.js 16 kết hợp TypeScript và Tailwind CSS để tối ưu hóa giao diện Single Page Application. Trợ lý ảo Catbot được vận hành dựa trên các mô hình ngôn ngữ lớn của OpenAI kết hợp với LangGraph để quản lý luồng hội thoại, lập kế hoạch tool-call và Cohere Rerank tùy chọn để tối ưu hóa kết quả tìm kiếm ngữ nghĩa.

Sau cùng, đóng góp chính và kết quả đạt được của đồ án tập trung vào ba giải pháp cốt lõi mang tính thực tiễn cao. Đóng góp thứ nhất là giải pháp Catbot RAG có planner tool-routing, có khả năng truy xuất hồ sơ sức khỏe thú cưng (Pet Profile), sản phẩm thật và kho tri thức để đưa ra các tư vấn cá nhân hóa và cảnh báo dị ứng khi dữ liệu đầu vào cho phép. Đóng góp thứ hai là thiết kế cơ chế chuyển giao linh hoạt giữa AI và con người (Human-in-the-loop), định tuyến từ Catbot sang nhân viên hỗ trợ khi người dùng yêu cầu gặp người thật, có dấu hiệu bức xúc hoặc câu hỏi vượt quá phạm vi tư vấn an toàn. Đóng góp thứ ba là triển khai quy trình giao dịch thương mại điện tử an toàn bằng khóa dòng dữ liệu bi quan (\texttt{SELECT FOR UPDATE}), khóa idempotency chống tạo đơn trùng và thanh toán QR qua SePay. Với đơn SePay, hệ thống tạo mã QR có thời hạn thanh toán 15 phút và reservation tồn kho có thêm khoảng grace 5 phút để worker có thể giải phóng hàng khi thanh toán quá hạn.

\section{Bố cục đồ án}
\label{section:1.4}

Phần còn lại của báo cáo đồ án tốt nghiệp này được tổ chức và trình bày theo cấu trúc gồm năm chương tiếp theo cùng hai phần phụ lục đính kèm, cụ thể như sau.

Chương 2 trình bày chi tiết về quá trình khảo sát hiện trạng thị trường và phân tích yêu cầu của hệ thống thương mại điện tử ThePawsome. Trong chương này, các yêu cầu chức năng của hệ thống được mô tả trực quan thông qua biểu đồ use case tổng quát và các biểu đồ use case phân rã cho từng nhóm tác nhân cụ thể. Đồng thời, quy trình nghiệp vụ mua hàng và xử lý thanh toán phức tạp được mô hình hóa bằng biểu đồ hoạt động, kết hợp với việc đặc tả chi tiết các use case nghiệp vụ cốt lõi và thiết lập các yêu cầu phi chức năng về hiệu năng vận hành và tiêu chuẩn an toàn thông tin.

Chương 3 giới thiệu chi tiết về nền tảng lý thuyết và công nghệ sử dụng trong đồ án. Nội dung chương tập trung phân tích cơ sở lý thuyết của kiến trúc Retrieval-Augmented Generation, kỹ thuật tìm kiếm ngữ nghĩa dựa trên cơ sở dữ liệu vector và phần mở rộng pgvector, thuật toán trộn điểm Reciprocal Rank Fusion cho mô hình tìm kiếm lai. Đồng thời, chương này cũng đưa ra các lập luận kỹ thuật giải thích lý do lựa chọn các công nghệ cốt lõi bao gồm FastAPI, Next.js, Redis, SePay, và các API trí tuệ nhân tạo của OpenAI và Cohere.

Chương 4 tập trung vào phần phân tích thiết kế, triển khai thực tế và đánh giá hiệu năng hệ thống. Chương này mô tả thiết kế kiến trúc phần mềm phân tầng thông qua biểu đồ gói UML, thiết kế chi tiết cơ sở dữ liệu quan hệ với sơ đồ thực thể liên kết ERD gồm 28 bảng dữ liệu chuẩn hóa, và thiết kế chi tiết các lớp đối tượng tham gia xử lý các nghiệp vụ quan trọng bằng biểu đồ trình tự. Kết quả xây dựng ứng dụng cũng được thống kê chi tiết về số lượng dòng code, cấu trúc mã nguồn và đánh giá thực nghiệm hiệu năng phản hồi của hệ thống.

Chương 5 trình bày sâu về các giải pháp kỹ thuật và đóng góp nổi bật của đồ án. Chương này phân tích chi tiết nguyên lý thiết kế và thuật toán triển khai của các giải pháp cốt lõi bao gồm: kiến trúc Catbot RAG tích hợp planner tool-routing và Pet Profile để tư vấn cá nhân hóa, cơ chế tìm kiếm lai kết hợp RRF và rerank tùy chọn, cơ chế giữ kho an toàn bằng khóa dòng dữ liệu bi quan, giải pháp thanh toán QR SePay kết hợp webhook, và quy trình chuyển giao hội thoại từ AI sang nhân viên hỗ trợ.

Chương 6 đưa ra các kết luận tổng quan và định hướng phát triển của đề tài trong tương lai. Chương này đánh giá khách quan các kết quả đạt được của sản phẩm so với mục tiêu ban đầu và so với các giải pháp hiện có trên thị trường, thẳng thắn nhìn nhận những mặt còn hạn chế, rút ra những bài học kinh nghiệm trong quá trình thiết kế hệ thống và đề xuất lộ trình nâng cấp các tính năng thông minh trong các giai đoạn tiếp theo.

Cuối cùng, phần phụ lục của báo cáo bao gồm hai phần nội dung bổ trợ. Phụ lục A cung cấp các hướng dẫn chi tiết về việc chuẩn bị môi trường, cài đặt, triển khai và kiểm thử mã nguồn của cả hai phân hệ Frontend và Backend ở môi trường local cũng như production. Phụ lục B trình bày các bảng đặc tả chi tiết cho các use case bổ sung của hệ thống nhằm đảm bảo tính đầy đủ và minh bạch cho tài liệu thiết kế.

\end{document}
